\section{Esperimenti}

\subsection{Esperimento preliminare}

E' stato eseguito un primo esperimento intermedio per verificare che la pipeline fosse corretta e che la rete fosse effettivamente in grado di apprendere la funzione di pricing.
Il comando usato e':

\begin{lstlisting}[style=pythonstyle, caption={Comando per un esperimento preliminare.}]
python scripts/run_experiment.py \
  --n-samples 10000 \
  --max-epochs 40 \
  --batch-size 1024 \
  --mc-n-paths 5000 \
  --mc-evaluation-samples 128
\end{lstlisting}

Questo run non rappresenta necessariamente l'esperimento finale del progetto, ma serve come controllo iniziale della correttezza dell'implementazione.
Gli esperimenti finali potranno usare piu' osservazioni, piu' epoche e un numero maggiore di path Monte Carlo.

\subsection{Visualizzazioni previste}

Le visualizzazioni principali sono:

\begin{itemize}
    \item andamento della training loss e validation loss;
    \item confronto tra prezzo Black-Scholes e prezzo predetto dalla rete;
    \item distribuzione degli errori;
    \item errore assoluto rispetto alla moneyness $S_0/K$;
    \item errore assoluto rispetto alla maturity $T$;
    \item errore assoluto rispetto alla volatilita' $\sigma$.
\end{itemize}

